% !TeX root = ../偏微分方程笔记.tex
% 由主文档通过 \input 载入；请编译主文档。
\chapter{绪论与阅读说明}
\label{chap:roadmap}

\section{研究对象与全书定位}

\emph{偏微分方程（partial differential equation，PDE）}研究多元未知函数与其偏导数之间的关系。温度如何扩散、波动如何传播、静电势如何由电荷和边界决定，都可以导向偏微分方程。学习这门课程既要理解方程如何产生，也要研究给定数据能否确定解，以及解具有怎样的性质。

本书按本科初学者的知识起点组织，以标量方程为主，从数学物理模型和经典求解方法出发，逐步过渡到弱解与现代分析方法。全书分为基础概念、经典方法与方程、弱解理论、进阶专题四部分。基础部分可作为第一门偏微分方程课程；后两部分供继续学习时使用，不预设读者在入门阶段已经掌握泛函分析。

正文围绕定义、推导、核心结论和完整例题展开。基础方法给出可核查的关键步骤；涉及较深实分析或泛函分析的结果，明确说明假设，并标注证明要点或进一步阅读出处。各章习题用于检验计算、证明和建模能力，可与正文例题配合使用。

初学时最容易遇到的困难，是能认出一个公式，却不知道为什么选用它。本书因此反复使用同一条思路：先认清未知量及数据，再由方程结构选择方法，最后把求出的表达式放回原问题中检查。例如，一根杆上的温度需要初温与端点热交换规律；一根弦的运动还需要初速度。两者都能展开成正弦函数，但每个正弦函数的振幅，一个随时间衰减，一个随时间振荡。理解这种差别，比记住两张互不相干的公式表更重要。

基础课程中的“求解”也有不同层次。算出表达式以后，需要解释它的收敛意义；证明解唯一以后，还要区分是否已经证明存在；得到能量递减以后，还要检查能否推出趋于零。正文中的例题尽量完成“条件、思路、计算、核查、解释”这一闭环。附录\ref{app:worked-exercises}提供各章代表习题的解答与提示，适合先独立尝试、再对照关键步骤使用。

\section{先修知识与学习目标}

基础部分需要多元微积分、常微分方程与线性代数：偏导数与链式法则、重积分与分部积分、散度定理、二阶线性常微分方程，以及实对称矩阵的特征值与正交对角化。傅里叶级数在第\ref{chap:fourier-series}章系统补充。进入第\ref{chap:sobolev}章前，应补足勒贝格积分、赋范空间与希尔伯特空间的基本概念；相关依赖列于附录\ref{app:analysis}和附录\ref{app:functional}。

完成基础部分后，应能够从守恒关系建立简单模型，正确提出定解问题，选用特征线、分离变量、积分变换或格林函数方法，并用极值原理与能量法讨论唯一性和稳定性。完成进阶部分后，应能够写出弱形式，检查存在唯一性定理的假设，理解正则性、非线性与数值近似各自带来的问题。

\section{全书课程大纲}

下列四个模块给出全书顺序，具体知识依赖在阅读路线中说明。

\begin{description}
  \item[基础概念与模型]
  第\ref{chap:concepts}章：偏微分方程的基本概念；
  第\ref{chap:modeling}章：物理模型与定解问题；
  第\ref{chap:characteristics}章：一阶方程与特征线法；
  第\ref{chap:classification}章：二阶线性方程的分类与标准形。
  \item[经典方法与三类基本方程]
  第\ref{chap:fourier-series}章：傅里叶级数与施图姆--刘维尔理论；
  第\ref{chap:separation}章：分离变量法与本征函数展开；
  第\ref{chap:wave-one}章：一维波动方程；
  第\ref{chap:wave-multi}章：多维波动方程与能量法；
  第\ref{chap:heat}章：热方程与扩散；
  第\ref{chap:harmonic}章：拉普拉斯方程与调和函数；
  第\ref{chap:green}章：泊松方程、格林函数与位势；
  第\ref{chap:transforms}章：积分变换与基本解。
  \item[弱解与现代分析基础]
  第\ref{chap:sobolev}章：分布、弱导数与索伯列夫空间；
  第\ref{chap:elliptic-weak}章：椭圆方程的弱解与变分方法；
  第\ref{chap:evolution-weak}章：演化方程的弱解与能量估计。
  \item[进阶专题与综合应用]
  第\ref{chap:conservation}章：非线性守恒律与熵解；
  第\ref{chap:nonlinear}章：非线性方程的基本方法；
  第\ref{chap:numerics}章：数值方法与综合建模。
\end{description}

附录\ref{app:calculus}整理微积分与常微分方程工具；附录\ref{app:analysis}整理积分、收敛与不等式；附录\ref{app:functional}整理泛函分析工具；附录\ref{app:worked-exercises}提供各章代表习题的解答与提示。前三个附录用于查阅和补课，不要求在阅读第一章前全部学完。

\section{阅读顺序与阶段性检查}

\begin{enumerate}
  \item \textbf{建立问题意识。}先读第\ref{chap:concepts}--\ref{chap:classification}章，能够区分方程、附加条件、解的类别和适定性问题，并独立求解简单输运方程。
  \item \textbf{掌握经典主线。}接着读第\ref{chap:fourier-series}--\ref{chap:green}章，比较波动、扩散与平衡三种机制。第\ref{chap:wave-multi}章的高维表示公式及第\ref{chap:separation}章的特殊函数可在第一遍阅读时选读。
  \item \textbf{统一求解方法。}第\ref{chap:transforms}章将频率分解、卷积与基本解联系起来。傅里叶变换也可在学完第\ref{chap:fourier-series}章后提前阅读，用于辅助热方程的推导。
  \item \textbf{进入弱解理论。}补足附录中的实分析与泛函分析工具，再读第\ref{chap:sobolev}--\ref{chap:evolution-weak}章，练习从经典方程到弱形式、从先验估计到存在性证明的转换。
  \item \textbf{按兴趣拓展。}第\ref{chap:conservation}章承接特征线与分布；第\ref{chap:nonlinear}章承接变分和演化方程；第\ref{chap:numerics}章中的差分部分可在基础课程后阅读，有限元部分则依赖弱形式。
\end{enumerate}

本科基础路线以能够独立提出问题、推导公式并核查解为标准；进阶路线以能够说明函数空间、估计与极限过程为标准，可在完成基础课程后继续阅读。

\section{统一符号与叙述约定}

空间维数记为正整数 $n$，空间点记为 $x=(x_1,\ldots,x_n)\in\mathbb R^n$，时间记为 $t$，终止时间记为 $T>0$。区域 $\Omega\subset\mathbb R^n$ 指非空连通开集，其边界和闭包分别为 $\partial\Omega$ 与 $\overline\Omega$；需要边界积分、法向量或迹时，另行说明边界的光滑性。外单位法向量记为 $\nu$。

标量未知函数记为 $u=u(x)$ 或 $u=u(x,t)$；$u_t$ 表示时间偏导数。梯度（gradient）、散度（divergence）和拉普拉斯算子（Laplacian）分别记为
\[
  \nabla u=(\partial_{x_1}u,\ldots,\partial_{x_n}u),\qquad
  \operatorname{div}F=\sum_{j=1}^n\partial_{x_j}F_j,\qquad
  \Delta u=\sum_{j=1}^n\partial_{x_jx_j}u,
\]
其中 $F=(F_1,\ldots,F_n)$ 为向量场，$\partial_\nu u=\nabla u\cdot\nu$ 为外法向导数。扩散系数记为 $\kappa>0$，波速记为 $c>0$，区间长度记为 $\ell>0$。源项通常记为 $f$；$g$、$h$ 表示具体问题中说明的已知初始或边界数据，各章均重新明确其含义。

$C^k$ 表示具有 $k$ 阶连续偏导数的函数类；$L^p$ 表示按几乎处处相等识别的勒贝格可积函数空间，其精确定义在附录\ref{app:analysis}补充。$\|v\|_{L^2(\Omega)}^2=\int_\Omega |v|^2\,dx$。多重指标、分布与索伯列夫空间在第\ref{chap:sobolev}章定义。常数的取值和依赖关系在具体估计中说明，不把不同上下文中的常数视为同一个数。

\section{参考资料与核对依据}

基础课程选题可对照 MIT 的公开课程大纲\cite{mit152}；弱导数、椭圆弱解和演化方程可配合 Hunter 的讲义\cite{hunter}阅读；分析与数值方法的衔接可对照 MIT 的相关课程\cite{mit303}。这些资料用于核对课程覆盖和继续阅读，本书的章节安排与例题设计按上述学习路线组织。

